%% Document généré pour rendu atomique — ne pas éditer à la main
%% Atome : fiche N12/S05 — 68a54b3cd64f4997b30c605caedb0674
%% Préambule sur mesure : 53 définitions inlinées, 24 paquets externes

\documentclass[11pt,a4paper]{article}

%% ── Paquets externes (noyau pour atome) ───────────────
\usepackage[utf8]{inputenc}
\usepackage[T1]{fontenc}
\usepackage{lmodern}
\usepackage[french]{babel}
\usepackage{geometry}
\usepackage{amsmath}
\usepackage{amssymb}
\usepackage{etoolbox}
\usepackage{xkeyval}
\usepackage{xifthen}
\usepackage{environ}
\usepackage{forloop}
\usepackage{xcolor}
\usepackage{ninecolors}
\usepackage{multirow}
\usepackage{colortbl}
\usepackage{tabularray}
\usepackage{multicol}
\usepackage[inline]{enumitem}
\usepackage{graphicx}
\usepackage{adjustbox}
\usepackage{tikz}
\usepackage[theorems,breakable,skins]{tcolorbox}
\usepackage{pifont}

%% ── Marges (alignées sur les livrets) ───────────────
\newgeometry{left=35pt,right=35pt,top=65pt,bottom=70pt}

%% ── Initialisations de bibliothèques ─────────────────
\usetikzlibrary{babel,shapes.geometric,arrows.meta,positioning,calc,matrix,fit,intersections}
\UseTblrLibrary{booktabs}
\UseTblrLibrary{varwidth}

%% ── Définitions du paquet seqenseigne (inlinées) ─────
%%    Fermeture transitive des macros effectivement utilisées
%%    par l'atome — au lieu de \usepackage{seqenseigne}
%%    qui chargerait ~110 paquets.
\makeatletter
\newcommand{\seqCouleurRougeVif}{red!85!black}
\newcommand{\seqCouleurRougeModere}{red!50!white}
\newcommand{\seqCouleurRougeMoyen}{red!25!white}
\newcommand{\seqCouleurRougePale}{red!10!white}
\newcommand{\seqCouleurJauneModere}{yellow!50!white}
\newcommand{\seqCouleurJaunePale}{yellow!10!white}
\newcommand{\seqCouleurBleuVif}{blue!85!black}
\newcommand{\seqCouleurBleuModere}{blue!50!white}
\newcommand{\seqCouleurBleuMoyen}{blue!25!white}
\newcommand{\seqCouleurBleuPale}{blue!10!white}
\newcommand{\seqCouleurCanardModere}{teal!50!white}
\newcommand{\seqCouleurCanardPale}{teal!10!white}
\newcommand{\seqCouleurMarronVif}{brown!85!black}
\newcommand{\seqCouleurMarronModere}{brown!50!white}
\newcommand{\seqCouleurMarronMoyen}{brown!25!white}
\newcommand{\seqCouleurMarronPale}{brown!10!white}
\newcommand{\seqCouleurOrangeModere}{orange!50!white}
\newcommand{\seqCouleurOrangePale}{orange!10!white}
\newcommand{\seqCouleurVioletVif}{violet!85!black}
\newcommand{\seqCouleurVioletMoyen}{violet!25!white}
\newcommand{\seqCouleurVioletModere}{violet!50!white}
\newcommand{\seqCouleurVioletPale}{violet!10!white}
\newcommand{\seqCouleurVertModere}{green!50!white}
\newcommand{\seqCouleurVertPale}{green!10!white}
\newcommand{\seqCouleurRoseVif}{magenta!85!black}
\newcommand{\seqCouleurRoseMoyen}{magenta!25!white}
\newcommand{\seqCouleurRoseModere}{magenta!50!white}
\newcommand{\seqCouleurRosePale}{magenta!10!white}
\newcommand{\seqCouleurLimeModere}{lime!50!white}
\newcommand{\seqCouleurLimePale}{lime!10!white}
\newcommand{\seqCouleurNoirVif}{black!85!black}
\newcommand{\seqCouleurNoirMoyen}{black!25!white}
\newcommand{\seqCouleurNoirModere}{black!50!white}
\newcommand{\seqCouleurNoirPale}{black!10!white}
\newcommand{\seqCouleurGrisModere}{gray!50!white}
\newcommand{\seqCouleurGrisPale}{gray!10!white}
\newcommand{\seq@setColors}[1]{
\ifthenelse{ \equal{#1}{nombres} }
{
\def\seq@couleuri@vif{\seqCouleurRougeVif}
\def\seq@couleuri@moyen{\seqCouleurRougeMoyen}
\def\seq@couleuri@modere{\seqCouleurRougeModere}
\def\seq@couleuri@pale{\seqCouleurRougePale}
\def\seq@couleurii@modere{\seqCouleurJauneModere}
\def\seq@couleurii@pale{\seqCouleurJaunePale}
}
{
\ifthenelse{ \equal{#1}{donnees} }
{
\def\seq@couleuri@vif{\seqCouleurBleuVif}
\def\seq@couleuri@moyen{\seqCouleurBleuMoyen}
\def\seq@couleuri@modere{\seqCouleurBleuModere}
\def\seq@couleuri@pale{\seqCouleurBleuPale}
\def\seq@couleurii@modere{\seqCouleurCanardModere}
\def\seq@couleurii@pale{\seqCouleurCanardPale}
}
{
\ifthenelse{ \equal{#1}{grandeurs} }
{
\def\seq@couleuri@vif{\seqCouleurMarronVif}
\def\seq@couleuri@moyen{\seqCouleurMarronMoyen}
\def\seq@couleuri@modere{\seqCouleurMarronModere}
\def\seq@couleuri@pale{\seqCouleurMarronPale}
\def\seq@couleurii@modere{\seqCouleurOrangeModere}
\def\seq@couleurii@pale{\seqCouleurOrangePale}
}
{
\ifthenelse{ \equal{#1}{geometrie} }
{
\def\seq@couleuri@vif{\seqCouleurVioletVif}
\def\seq@couleuri@moyen{\seqCouleurVioletMoyen}
\def\seq@couleuri@modere{\seqCouleurVioletModere}
\def\seq@couleuri@pale{\seqCouleurVioletPale}
\def\seq@couleurii@modere{\seqCouleurVertModere}
\def\seq@couleurii@pale{\seqCouleurVertPale}
}
{
\ifthenelse{ \equal{#1}{algorithmique} }
{
\def\seq@couleuri@vif{\seqCouleurRoseVif}
\def\seq@couleuri@moyen{\seqCouleurRoseMoyen}
\def\seq@couleuri@modere{\seqCouleurRoseModere}
\def\seq@couleuri@pale{\seqCouleurRosePale}
\def\seq@couleurii@modere{\seqCouleurLimeModere}
\def\seq@couleurii@pale{\seqCouleurLimePale}
}
{
\def\seq@couleuri@vif{\seqCouleurNoirVif}
\def\seq@couleuri@moyen{\seqCouleurNoirMoyen}
\def\seq@couleuri@modere{\seqCouleurNoirModere}
\def\seq@couleuri@pale{\seqCouleurNoirPale}
\def\seq@couleurii@modere{\seqCouleurGrisModere}
\def\seq@couleurii@pale{\seqCouleurGrisPale}
}
}
}
}
}
}
\newcommand{\acompleter}[1]{\;\underline{\phantom{#1}}\;}
\newtcolorbox{boiteTitreSection}{
colback=\seq@couleuri@pale,
colframe=\seq@couleuri@vif,
coltext=\seq@couleuri@vif
}
\define@cmdkey[seq]{boiteContenuFlashcard}{titre}[Exemples]{}
\newenvironment{seqBoiteContenuFlashcard}[1][]{
\setkeys[seq]{boiteContenuFlashcard}{titre,#1}
\begin{tcolorbox}[title=\cmdseq@boiteContenuFlashcard@titre
,natural height
,colback=\seq@couleurii@pale
,colframe=\seq@couleuri@vif
,colbacktitle=\seq@couleuri@vif
,fonttitle=\bfseries
,enhanced
,attach boxed title to top left={yshift=-5pt,xshift=20pt}
]
}{
\end{tcolorbox}
}
\define@cmdkey[seq]{boiteFillContenuFlashcard}{titre}[Exemples]{}
\newenvironment{seqBoiteFillContenuFlashcard}[1][]{
\setkeys[seq]{boiteFillContenuFlashcard}{titre,#1}
\begin{tcolorbox}[title=\cmdseq@boiteFillContenuFlashcard@titre
,height fill
,colback=\seq@couleurii@pale
,colframe=\seq@couleuri@vif
,colbacktitle=\seq@couleuri@vif
,fonttitle=\bfseries
,enhanced
,attach boxed title to top left={yshift=-5pt,xshift=20pt}
]
}{
\end{tcolorbox}
\newpage
}
\define@cmdkey[seq]{colItem}{nbCols}[2]{}
\define@cmdkey[seq]{colItem}{label}[$\bullet$]{}
\newenvironment{seqColItem}[1][]{
\setkeys[seq]{colItem}{nbCols,label,#1}
\ifthenelse{\cmdseq@colItem@nbCols > 1}
{\begin{multicols}{\cmdseq@colItem@nbCols}}{}
\begin{itemize}[label=\cmdseq@colItem@label]
}{
\end{itemize}
\ifthenelse{\cmdseq@colItem@nbCols > 1}{\end{multicols}}{}
}
\newcommand{\seqSetColorsTheme}[1]{
\seq@setColors{#1}
}
\newcommand{\seqSetCodeNiveau}[1]{\gdef\seq@code@niveau{#1}}
\newcommand{\seqSetCodeSequence}[1]{\gdef\seq@code@sequence{#1}}
\newcommand{\seqTitreSection}[1]{
\ifdocstd
\newpage
\newgeometry{left=35pt,right=35pt,top=65pt,bottom=70pt}
\fi
\begin{boiteTitreSection}
\begin{minipage}[c][20pt][c]{\linewidth}
\centering \Large \textbf{#1}
\end{minipage}
\end{boiteTitreSection}
\par\vspace{10pt}
}
\newcommand\seqCreeCompteurs{
\newcounter{NotionNum}
\newcounter{DefinitionNum}
\newcounter{ProprieteNum}
\newcounter{MethodeNum}
\newcounter{SerieExosNum}
\newcounter{ExoNum}
\newcounter{AnnexeNum}

\renewcommand{\theHExoNum}{\theSerieExosNum-\theExoNum}
\renewcommand{\theHNotionNum}{\theSerieExosNum-\theNotionNum}
\renewcommand{\theHMethodeNum}{\theSerieExosNum-\theMethodeNum}
\renewcommand{\theHDefinitionNum}{\theSerieExosNum-\theDefinitionNum}
\renewcommand{\theHProprieteNum}{\theSerieExosNum-\theProprieteNum}
}
\newcommand{\seqFrac}[2]{\dfrac{\strut{} #1}{\strut{} #2}}
\newif\ifdocstd
\makeatother

%% ── Macros transverses ───────────────────────────────
\providecommand{\acompleter}[1]{\;\underline{\phantom{#1}}\;}

%% Thème de la séquence
\seqSetColorsTheme{nombres}
\seqSetCodeNiveau{N12}
\seqSetCodeSequence{S05}

\begin{document}

%% Fallback pour les \theHxxx normalement définis par hyperref.
\providecommand{\theHExoNum}{\theExoNum}
\providecommand{\theHNotionNum}{\theNotionNum}
\providecommand{\theHMethodeNum}{\theMethodeNum}
\providecommand{\theHDefinitionNum}{\theDefinitionNum}
\providecommand{\theHProprieteNum}{\theProprieteNum}

\seqCreeCompteurs

\seqTitreSection{Modéliser et résoudre des problèmes faisant intervenir le calcul littéral: équation du premier degré; équation-produit se ramenant au premier degré; équation de la forme \mbox{$x^2 = a$}.}
\begin{seqBoiteContenuFlashcard}[titre=Méthode]
Modéliser un problème signifie \acompleter{traduire l'énoncé}, qui est écrit en langage \acompleter{naturel}, en langage \acompleter{mathématique}:
			\begin{seqColItem}[nbCols=1,label=\ding{224}]
				\item La valeur cherchée dans le problème devient \acompleter{l'inconnue de l'équation}.
				\item Le type d'équation qu'on va écrire dépend fortement \acompleter{de l'énoncé} et il n'y a pas de procédure précise à suivre pour construire l'équation. \newline Cependant, on se basera souvent sur \acompleter{la compréhension des mathématiques en langage naturel}, \acompleter{la connaissance des formules de calcul de grandeurs}, et éventuellement \acompleter{les formules scientifiques \og classiques \fg{}}.
			\end{seqColItem}
\end{seqBoiteContenuFlashcard}
\begin{seqBoiteContenuFlashcard}[titre=Méthode]
Une équation du premier degré est de la forme \og \acompleter{$ax+b=0$} \fg{} où $a$ et $b$ sont \acompleter{des nombres} (avec \acompleter{$a \neq 0$}) et $x$ est \acompleter{l'inconnue}. \newline Résoudre cette équation veut dire trouver par quelle \acompleter{valeur numérique} il faut remplacer \acompleter{$x$} pour que l'équation \acompleter{soit vraie}. \newline Pour ce faire, on a le droit d'effectuer les \acompleter{transformations} suivantes, qui ne changent pas \acompleter{les solutions de l'équation}:
			\begin{seqColItem}[nbCols=1,label=\ding{224}]
				\item Ajouter \acompleter{la même quantité} aux deux \acompleter{membres} de l'équation.
				\item Multiplier chaque \acompleter{membre} de l'équation par \acompleter{la même quantité non nulle}.
			\end{seqColItem}
\end{seqBoiteContenuFlashcard}
\begin{seqBoiteContenuFlashcard}[titre=Méthode]
En utilisant le symbole \og $\Leftrightarrow$ \fg{} qui signifie que \acompleter{l'équation précédente} a les mêmes solutions que \acompleter{l'équation suivante}, on a:
			\begin{align*}
				& ax+b &= \quad & 0 \\
				\Leftrightarrow \quad & ax+b\acompleter{-b} &= \quad & 0\acompleter{-b} & \text{(en \acompleter{ajoutant $(-b)$} aux 2\,membres)} \\
				\Leftrightarrow \quad & \acompleter{ax} &= \quad & \acompleter{-b} & \\
				\Leftrightarrow \quad & ax \acompleter{\div a} &= \quad & -b \acompleter{\div a} & \text{(en \acompleter{divisant par $a$} les 2\,membres)} \\
				\Leftrightarrow \quad & \acompleter{x} &= \quad & \acompleter{-\seqFrac{b}{a}}
			\end{align*}
			Finalement, l'ensemble solution est $\mathcal{S}=\{\acompleter{-\seqFrac{b}{a}}\}$.
\end{seqBoiteContenuFlashcard}
\begin{seqBoiteContenuFlashcard}[titre=Méthode]
Une équation-produit se ramenant au premier degré est de la forme \og \acompleter{$(ax+b)(cx+d)=0$} \fg{} où $a$, $b$, $c$ et $d$ sont \acompleter{des nombres} (avec \acompleter{$a \neq 0$} et \acompleter{$c \neq 0$}) et $x$ est \acompleter{l'inconnue}. \newline En utilisant ce qui précède, on a:
			\begin{align*}
				& (ax+b)(cx+d) &= \quad & 0 \\
				\Leftrightarrow \quad & \acompleter{ax+b=0} & \text{ou\quad} & \acompleter{cx+d=0} & \text{(un produit est nul quand \acompleter{un des facteurs l'est})} \\
				\Leftrightarrow \quad & x=\acompleter{-\seqFrac{b}{a}} & \text{ou\quad} & x=\acompleter{-\seqFrac{d}{c}} 
			\end{align*}
			Finalement, l'ensemble solution est $\mathcal{S}=\{\acompleter{-\seqFrac{b}{a}}; \acompleter{-\seqFrac{d}{c}}\}$.
\end{seqBoiteContenuFlashcard}
\begin{seqBoiteContenuFlashcard}[titre=Méthode]
Pour résoudre une équation de la forme $x^2=a$, on doit considérer \acompleter{2\,cas disjoints}:
			\begin{seqColItem}[nbCols=1,label=\ding{224}]
				\item Si $a$ est \acompleter{strictement négatif}, l'ensemble solution est \acompleter{vide} car un \acompleter{carré} est \acompleter{toujours positif ou nul}: $\mathcal{S}=\acompleter{\emptyset}$ 
				\item Si $a$ est \acompleter{positif ou nul}, on a:
					\begin{align*}
						& x^2 &= \quad & a \\
						\Leftrightarrow \quad & x^2\acompleter{-a} &= \quad & a\acompleter{-a} & \text{(en \acompleter{ajoutant $(-a)$} aux 2\,membres)} \\
						\Leftrightarrow \quad & \acompleter{x^2-a} &= \quad & \acompleter{0} & \\
						\Leftrightarrow \quad & x^2-\acompleter{(\sqrt{a})^2} &= \quad & 0 & \text{(car \acompleter{$a=(\sqrt{a})^2$})} \\
						\Leftrightarrow \quad & (\acompleter{x+\sqrt{a}})(\acompleter{x-\sqrt{a}}) &= \quad & 0 & \text{(en utilisant \acompleter{la 3e identité remarquable})} \\
						\Leftrightarrow \quad & x=\acompleter{-\sqrt{a}} & \text{ou\quad} & x=\acompleter{\sqrt{a}} 
					\end{align*}
					Finalement, l'ensemble solution est $\mathcal{S}=\{\acompleter{-\sqrt{a}}; \acompleter{\sqrt{a}}\}$.
			\end{seqColItem}
\end{seqBoiteContenuFlashcard}
\begin{seqBoiteFillContenuFlashcard}\end{seqBoiteFillContenuFlashcard}

\end{document}